\documentclass[12pt]{article}

\usepackage[a4paper, total={6in, 8in}]{geometry}
\usepackage{textcomp}

\begin{document}
\setlength{\parindent}{0pt}

\def \t {\textrightarrow}
\def \v {\vspace{1em}}
\def \bi {\begin{itemize}}
\def \ei {\end{itemize}}
\def \s[#1] {\section*{#1}}
\def \ss[#1] {\subsection*{#1}}
\def \sss[#1] {\subsubsection*{#1}}

\s[Costituzione]
La nostra costituzione è il frutto dell'assemblea costituente, eletta il 2 giugno 1946 \r redime testo definitivo nel dicembre 47 \\
Viene firmata dal capo provvisorio De Nicola

\v

Sostituisce lo statuto albertino, che era in vigore da 100 anni \t era stato concesso da Carlo Alberto, con le guerre di indipendenza \t è l unico statuto che non viene abrogato (mentre quelli concessi dagli altri si) \\
Statuto, non costituzione \t perche viene concesso unilateralmente dall'alto \t è un regolamento, delle norme dettate dal re senza una consultazione \\
È una costituzione ottriata = concessa dall'alto \t il sovrano declina una serie di norme \\
Statuto segna il passaggio da monarchia a monarchia costituzionale \\
Governo è responsabile davanti al re in monarchia, mentre al parlamento in monarchia costituzionale \\
Mussolini mantiene lo statuto perchè era una costituzione flessibile (mentre quella attuale è rigida) \\
Una costituzione flessibile puo essere modificata con maggioranza parlamentare semplicemente \t non richiede un iter parlamentare complesso \\
Mussolini quinid naviga nelle regole dello statuto

\v

Sulla scia di questa esperienza, la costituzione attuale è rigida: non puo essere cambiata con una legge ordinaria \t serve iter molto complesso \\
Nulla tranne l'articolo 7, che quando viene concepito la possibilita di modifica \t regola rapporti stato - chiesa, non è necessaria "procedimento di revisione costituzionale" \\
La costituzione del 48 comprende nel 7 i patti Lateranensi \t voluto fortemente da De Gasperi \t e Togliatti non lo osteggia, per evitare scontro \\
Nel 1984 Craxi, presidente del consiglio, firma con il vaticano nuovo concordato \t che ha fatto diventare lo stato italiano uno stato laico: religione cattolica riconosciuta come di maggioranza, ma non come religione di stato \\
Adesso non è + obbligatoria religione a scuola, inoltre il matrimonio religioso non ha immediato valore civile

\v

I primi 12 principi hanno il compito di riassumere cio che viene approfondito nei vari articoli successivi

\v

La nostra costituzione è antifascista \t ci sono degli aspetti che hanno radici nella nostra storia \t viene elaborata infatti nel 46 \\
Viene scritta con l'intento di evitare che si possa ripetere quell'esperienza \\
\bi
  \item è composta dai 12 principi fondamentali
  \item una prima parte che tratta dei diritti e doveri dei cittadini
  \item secoda parte che tratta della struttura dello stato
  \item epilogo
\ei
La costituzione afferma che la sovranita spetta al popolo \t ma se la sovranita è popolare, il popolo ha anche dei doveri nella gestione del potere \\
La costituzione garantisce dei diritti, ma il popolo ha anche doveri \\
Per esempio articolo 4: la repubblica riconosce il diritto al lavoro \t ma poi, ogni cittadino ha il dovere di partecipare a un attivita \\
Lo stato da possibilita, e il cittadino si impegna ad accogliere queste opportunita \\
La costituzione non è un insieme di leggi, non è un codice \t ma è un insieme di norme: indica i criteri che devono ispirare la giurisprudenza \t il codice civile e penale devono rispettare i principi fondamentali dettati dall costituzione
\end{document}
